\documentclass[10pt,conference,compsocconf]{IEEEtran}

\usepackage{hyperref}
\usepackage{graphicx}
\usepackage{amsmath}
\usepackage{amssymb}
\usepackage{booktabs}
\usepackage{url}

\begin{document}

\title{Predicting coronary heart disease from a 300\,000-respondent health
survey: where the accuracy actually comes from}

% TEAM: replace with the three full names before submitting.
\author{
  \IEEEauthorblockN{Alban --- teammate 2 --- teammate 3}
  \IEEEauthorblockA{\textit{CS-433 Machine Learning, EPFL}}
}

\maketitle

\begin{abstract}
We predict coronary heart disease from the 2015 BRFSS telephone survey
(328\,135 respondents, 321 columns, 8.8\,\% positives) using only NumPy. Our
best model, $L_2$-regularised logistic regression fitted by Newton's method,
reaches an out-of-fold $F_1$ of $0.441$. Our ablation shows that the decision
threshold, not the model or the cleaning, carries almost all of the
performance: tuning it raises $F_1$ from $0.259$ to $0.441$, while our entire
hand-crafted decoding of the survey codes is worth less than $0.005$. We also
find the one-hot design matrix to be rank-deficient, so the normal equations of
least squares have no unique solution until redundant columns are removed.
\end{abstract}

\section{Introduction}

The Behavioral Risk Factor Surveillance System (BRFSS) is an annual telephone
survey of U.S. residents. Given the 321 answers of a respondent, the task is to
predict whether that person has been diagnosed with myocardial infarction or
coronary heart disease (MICHD). Two properties of the data drive every decision
in this report. First, the classes are strongly imbalanced: only $8.8\,\%$ of
the respondents are positive, so accuracy is meaningless (predicting
``healthy'' for everybody already scores $91.2\,\%$) and the competition is
scored with the $F_1$ score of the positive class. Second, the data is a
questionnaire, not a set of measurements: $44.8\,\%$ of the entries are
missing, and the numbers stored in the file are mostly category codes.

We built a pipeline that turns the survey codes into meaningful numbers, fits a
regularised logistic model, and treats the decision threshold as a
hyperparameter. All results below are 5-fold cross-validated, with the
preprocessing re-estimated inside every training fold so that no statistic from
a validation fold reaches the model.

\section{Data and feature engineering}

\subsection{What the raw columns look like}

Our exploratory analysis (\texttt{eda.py}) shows that $220$ columns take at
most ten distinct values and are categorical codes, $80$ are genuinely
numerical, and the rest are administrative (identifiers, interview dates,
survey design weights). Missingness is very uneven --- $147$ columns are
missing more than half of the time and $99$ more than $90\,\%$ --- because
whole blocks of questions are only asked to a sub-sample of respondents.

\subsection{Cleaning}

The survey reserves codes for non-answers ($7$/$9$, $77$/$99$, $777$/$999$
depending on the field width) and uses compound codes for quantities: $1xx$
means $xx$ times per week, $2xx$ per month and $3xx$ per year; weight is either
in pounds or in kilograms offset by $9000$; height is either \texttt{ftii} or
centimetres offset by $9000$. Taken literally, a respondent who refuses to
answer a frequency question appears to exercise $999$ times a week. We decode
the compound codes into a single unit and map the reserved codes to missing
values.

Columns with few distinct values are one-hot encoded against their most
frequent level; the others are winsorised at the 1st and 99th percentiles,
median-imputed and standardised. For every column missing more than $2\,\%$ of
the time we add a ``was missing'' indicator: in a survey a non-answer is itself
an answer. Finally we append the square of every numerical feature, so the
linear model can represent the U-shaped effect of age or body mass index.

\subsection{A rank-deficient design matrix}

One-hot encoding a survey creates exact linear dependencies: a missingness
indicator is often already implied by the level dummies of the same column. On
our default pipeline the design matrix has $415$ columns but $X^\top X$ has
rank $344$, so \texttt{least\_squares} raises a singular-matrix error. Either
one regularises (any $\lambda>0$ makes $X^\top X + 2N\lambda I$ invertible) or
one removes the redundant columns. We do the latter, with an incremental
Cholesky factorisation of the Gram matrix that keeps a column only when it adds
a direction the selected ones cannot represent. This restores full rank without
changing the span, and lets all six methods run on identical features.

\section{Models}

\subsection{The six required methods}

We implemented the six methods of the project statement in
\texttt{implementations.py}, solving every linear system with
\texttt{np.linalg.solve} rather than inverting a matrix. For the iterative
methods the step size is not a free choice: gradient descent on the mean
squared error converges if and only if $\gamma < 2/\lambda_{\max}(X^\top X/N)$.
Measuring $\lambda_{\max} = 170.2$ on our design matrix predicts a stability
limit of $\gamma = 0.0117$, which the data confirms exactly: $\gamma = 0.01$
converges to a loss of $0.137$ while $\gamma = 0.012$ diverges to $10^{10}$
within 300 iterations. For stochastic gradient descent the relevant bound
involves the curvature of a \emph{single} sample,
$2/\max_n \lVert x_n \rVert^2 = 1.45 \times 10^{-4}$, two orders of magnitude
smaller, which is why the two methods appear with very different step sizes in
Table~\ref{tab:methods}. Adding the quadratic features multiplies
$\lambda_{\max}$ by $15$, so a step size tuned before the feature expansion
silently diverges afterwards --- a trap we fell into.

\subsection{Newton's method}

Gradient descent on this data set is slow because the features have very
different curvatures. Since the design matrix has only a few hundred columns,
we can afford the exact Newton step
\[
  w \leftarrow w - H^{-1} \nabla \mathcal{L}(w),
  \qquad
  H = \tfrac{1}{N} X^\top S X + 2\lambda I,
\]
where $S = \mathrm{diag}\left(\sigma(1-\sigma)\right)$ collects the curvature
of each sample. This is one $(P,P)$ linear solve per iteration on top of the same matrix
products gradient descent already performs. Newton converges in about ten
iterations where 300 gradient steps are still far from the optimum, for a
comparable total cost (Table~\ref{tab:methods}).

\subsection{The decision threshold}

A model trained on a set with $8.8\,\%$ positives puts almost every predicted
probability below $0.5$, so it flags far too few respondents: with the default
threshold our model reaches $F_1 = 0.259$ despite an accuracy of $91.7\,\%$,
its highest accuracy of all. We therefore treat the threshold as a hyperparameter
and select the value maximising $F_1$ on out-of-fold predictions. The whole
$F_1$ curve costs a single sort, since after ordering the scores the number of
true positives above a threshold is a cumulative sum.

\section{Results and discussion}

\begin{table}[t]
\centering
\caption{The six required methods and our variant, 5-fold cross-validation on
the same features, with the decision threshold tuned on the training fold.}
\label{tab:methods}
\small
\begin{tabular}{lccc}
\toprule
Method & $F_1$ & Accuracy & Time \\
\midrule
% BEGIN METHODS
Least squares (normal eq.) & $0.430 \pm 0.006$ & $0.876$ & $55$\,s \\
Ridge regression & $0.430 \pm 0.006$ & $0.875$ & $71$\,s \\
MSE gradient descent & $0.397 \pm 0.004$ & $0.855$ & $166$\,s \\
MSE stochastic GD & $0.385 \pm 0.012$ & $0.851$ & $62$\,s \\
Logistic regression (GD) & $0.405 \pm 0.005$ & $0.866$ & $174$\,s \\
Reg. logistic (GD) & $0.404 \pm 0.005$ & $0.866$ & $165$\,s \\
Reg. logistic (Newton) & $0.441 \pm 0.006$ & $0.877$ & $191$\,s \\
Newton, default threshold & $0.259 \pm 0.003$ & $0.917$ & $196$\,s \\
% END METHODS
\bottomrule
\end{tabular}
\end{table}

\begin{table}[t]
\centering
\caption{Ablation of the pipeline, 3-fold cross-validation, regularised
logistic regression fitted by Newton's method.}
\label{tab:ablation}
\small
\begin{tabular}{lcc}
\toprule
Configuration & $F_1$ & $\Delta F_1$ \\
\midrule
% BEGIN ABLATION
Full pipeline & $0.440 \pm 0.002$ & --- \\
$-$ quadratic terms & $0.439 \pm 0.001$ & $-0.000$ \\
$-$ one-hot encoding & $0.439 \pm 0.001$ & $-0.001$ \\
$-$ missingness indicators & $0.433 \pm 0.001$ & $-0.007$ \\
$-$ reserved-code masking & $0.438 \pm 0.002$ & $-0.002$ \\
$-$ compound decoding & $0.439 \pm 0.001$ & $-0.001$ \\
$-$ columns $>50\%$ missing & $0.426 \pm 0.001$ & $-0.013$ \\
No cleaning at all & $0.430 \pm 0.002$ & $-0.010$ \\
$+$ balanced class weights & $0.438 \pm 0.002$ & $-0.002$ \\
% END ABLATION
\bottomrule
\end{tabular}
\end{table}

\paragraph{The threshold dominates everything.} Comparing
Table~\ref{tab:methods} with the same model evaluated at the default threshold
shows the single largest effect we measured: $F_1$ falls from $0.441$ to
$0.259$ while accuracy \emph{rises} from $0.877$ to $0.917$. The two metrics
disagree because they answer different questions: optimising the wrong one
produces a model that misses most patients in order to protect a majority it
was never asked about. Any report on this data set that quotes accuracy is
reporting the prevalence of the disease, not the quality of the model.

\paragraph{Our careful cleaning was largely wasted.} Removing the reserved-code
masking, the compound decoding, the one-hot encoding or the quadratic terms
costs at most $0.002$ of $F_1$ each, and a deliberately naive pipeline --- every
column read as a raw number, median-imputed and standardised --- lands $0.010$
below the full pipeline, less than the gap between gradient descent and Newton.
Two explanations fit: the codes $7$ and $9$ are numerically far from the valid
range $1$--$5$, so a linear model isolates them with one coefficient anyway,
and the BRFSS derives dozens of variables from the same raw answers, so
corrupting one representation of age leaves several clean ones. We report this
negative result because it would have changed how we spent our time.

\paragraph{What did matter} concerns \emph{missingness} and
\emph{optimisation}, not the content of the answers. Keeping every column,
including those missing more than $90\,\%$ of the time, is worth $+0.013$ over
dropping everything above $50\,\%$ missingness, and the missingness indicators
are worth $+0.007$ --- by far the largest contribution of the pipeline. A
question asked only to a sub-sample identifies that sub-sample: the pattern of
non-answers encodes which survey module a respondent was routed through.
Replacing gradient descent by Newton's method is worth $+0.035$ at comparable
cost, because $300$ gradient steps are nowhere near the optimum. And tuning the
threshold is worth $+0.18$.

\paragraph{Limitations.} Our model is linear; the quadratic terms are a crude
substitute for the interactions a tree ensemble would find and bought only
$+0.0004$, while pairwise products of the twenty most correlated features added
$+0.0003$ for $145$ columns, so we dropped them. Balanced class weights made
things slightly worse once the threshold was tuned, as expected: reweighting
classes and moving the threshold do the same thing.

\section{Conclusion}

Regularised logistic regression fitted by Newton's method on the encoded BRFSS
survey reaches an out-of-fold $F_1$ of $0.441$ at $87.7\,\%$ accuracy, using
only NumPy. The ablation ranks our decisions unambiguously: the threshold
($+0.18$) is worth five times the optimiser ($+0.035$), itself worth three
times the feature pipeline ($+0.010$). On an imbalanced problem, what you
optimise matters more than how.

\end{document}
